\section{Introduction}
\label{sec:introduction}

The political feasibility of the \dia{} depends partly on whether
skeptics can credibly argue that algorithm parameters can be tuned
to produce partisan outcomes.
A hostile legislator or redistricting commission might argue: ``We are
just setting technical parameters.
Our choices of population tolerance, county preservation weight, and
convergence threshold are standard algorithmic practice.
The partisan outcome of our parameter choices is coincidental.''

U.2~\citep{b17paper} answered the descriptive version of this question:
what is the observed variation in partisan outcomes as parameters vary?
The answer was reassuring: at most 0.3 seats nationally across the full
parameter range.
But U.2 did not frame the question adversarially.
It asked ``what happens when we vary parameters one at a time?'' not
``what is the maximum a hostile actor can achieve by jointly optimizing
all parameters for partisan advantage?''

This paper addresses the adversarial framing directly.

\subsection{Adversarial Threat Model}

We model a hostile redistricting administrator who:
\begin{enumerate}
\item Knows the full \bisect{} algorithm specification;
\item Has freedom to set any parameter within its admissible range;
\item Can run the algorithm with different parameter settings before
  Census data are released (to understand the distribution of partisan
  outcomes without actually observing them), or after release
  (if pre-registration is not required);
\item Will choose the parameter setting that maximizes partisan advantage
  for their preferred party.
\end{enumerate}

The administrator is constrained only by the requirement that parameters
fall within their technically-justified ranges.
They cannot set $\ufactor = 50\%$ (which would produce grossly
unequal population districts) or $\acounty = 1{,}000$ (which would
effectively prevent any county splits, violating the one-person-one-vote
requirement in multi-county districts).

\subsection{Main Contributions}

\begin{enumerate}
\item We define the \emph{most partisan parameter configurations}
  (one Republican-favorable, one Democratic-favorable) for the \bisect{}
  five-parameter space (Section~\ref{sec:most-partisan}).
\item We report the end-to-end partisan outcome difference between the two
  configurations across all 50 states (Section~\ref{sec:results}).
\item We explain \emph{why} partisanship is parameter-insensitive by
  decomposing the mechanism at the census-tract level
  (Section~\ref{sec:adversarial}).
\item We describe the \dia{} pre-registration defense and show it
  eliminates the gaming vector regardless of the 0.3-seat bound
  (Section~\ref{sec:preregistration}).
\end{enumerate}

\subsection{Relationship to U.2}

This paper uses U.2's 50-state dataset as its empirical foundation
but reorganises the analysis around the adversarial question.
U.2's primary finding---``parameter tuning produces $\leq 0.3$ seats of
variation nationally''---becomes this paper's \emph{answer} to the
question: ``what is the maximum a hostile actor can achieve?''
We add the adversarial interpretation, the worst-case configuration
analysis, and the DIA defense.

\subsection{Road Map}

Section~\ref{sec:parameters} reviews the five parameters and their
admissible ranges.
Section~\ref{sec:adversarial} frames the reanalysis adversarially.
Section~\ref{sec:most-partisan} defines the most partisan parameter
configurations.
Section~\ref{sec:results} reports the 50-state outcome comparison.
Section~\ref{sec:preregistration} describes the DIA defense.
Section~\ref{sec:conclusion} concludes.
